\subsection{Вычисление оценки МНК}

Вычислим матрицу плана \( \Phi \) для выборки при помощи скрипта на Python с~применением библиотеки NumPy для удобства матричных вычислений:

\begin{lstlisting}[caption=Расчёт матрицы плана]
import numpy as np

X = np.array(task["x"])

def phi(x):
  return [np.ones_like(x), x, x**2]

Phi = np.vstack(phi(X))

print(*Phi, sep='\n')
\end{lstlisting}

Итого имеем:

\begin{table}[H]
\centering
\begin{tabular}{|c|c|c|c||c|c|c|c|}
\hline
$i$ & $\varphi_0(X_i)$ & $\varphi_1(X_i)$ & $\varphi_2(X_i)$ & $i$ & $\varphi_0(X_i)$ & $\varphi_1(X_i)$ & $\varphi_2(X_i)$ \\ \hline\hline
1 & 1.0 & 1.0 & 1.0000 & 11 & 1.0 & 2.0 & 4.0000 \\ \hline
2 & 1.0 & 1.1 & 1.2100 & 12 & 1.0 & 2.1 & 4.4100 \\ \hline
3 & 1.0 & 1.2 & 1.4400 & 13 & 1.0 & 2.2 & 4.8400 \\ \hline
4 & 1.0 & 1.3 & 1.6900 & 14 & 1.0 & 2.3 & 5.2900 \\ \hline
5 & 1.0 & 1.4 & 1.9600 & 15 & 1.0 & 2.4 & 5.7600 \\ \hline
6 & 1.0 & 1.5 & 2.2500 & 16 & 1.0 & 2.5 & 6.2500 \\ \hline
7 & 1.0 & 1.6 & 2.5600 & 17 & 1.0 & 2.6 & 6.7600 \\ \hline
8 & 1.0 & 1.7 & 2.8900 & 18 & 1.0 & 2.7 & 7.2900 \\ \hline
9 & 1.0 & 1.8 & 3.2400 & 19 & 1.0 & 2.8 & 7.8400 \\ \hline
10 & 1.0 & 1.9 & 3.6100 & 20 & 1.0 & 2.9 & 8.4100 \\ \hline
\end{tabular}
\caption{Матрица плана $\Phi$ в табличном виде}
\end{table}

Заметим, что \( \textrm{rank}\,\Phi=3 \), поэтому можем применить метод наименьших квадратов для поиска оценки неизвестных параметров \( \hat{\vec{\omega}} \) по формуле \eqref{eq:approximation}.

Воспользуемся следующим скриптом на Python для вычисления оценки: 

\begin{lstlisting}[caption=Расчёт ОМНК]
Z = np.array(task["t"])

omega = np.linalg.inv(Phi @ Phi.T) @ Phi @ Z

print(omega)
\end{lstlisting}

Итого имеем:
\[
    \hat{\vec{\omega}}=\begin{pmatrix}
        4.94827296 & -3.93485988 & 0.98541809
    \end{pmatrix}^T
\]

Таким образом, искомая линия регрессии имеет вид:
\[
f(x) \approx 4.948 - 3.935x + 0.985x^2
\]